\documentclass[11pt]{article}
\usepackage{amsmath, amssymb, amsthm}
\usepackage{hyperref}
\usepackage{bm}

% Theorem environments
\newtheorem{definition}{Definition}
\newtheorem{theorem}{Theorem}
\newtheorem{proposition}{Proposition}
\newtheorem{corollary}{Corollary}
\newtheorem{example}{Example}


\title{Ffix0 as a Structural Fixed Point in Multi-Model Reasoning Systems}
\author{Masaomi Chiba}
\date{2025-12-01}

\begin{document}
\maketitle

\begin{abstract}
Large Language Models (LLMs) exhibit distinct internal reasoning geometries depending on their architectural principles. 
When multiple LLMs interact indirectly through human-mediated dialogue, their reasoning structures may be interpreted in terms of \emph{content expansion} versus \emph{structural minimization}. 
This paper formalizes the notion of \textbf{$\mathrm{Ffix0}$}, originally emerging from such meta-dialogical observations, as a canonical fixed point of structural coherence minimizing a generalized notion of structural entropy, or ``Fuss''. 
We demonstrate that $\mathrm{Ffix0}$ naturally arises as the terminal object in a category of reasoning functors and that it provides a unifying notion of coherence across heterogeneous model architectures.
\end{abstract}

\section{Introduction}

Recent meta-dialogical experiments involving multiple state-of-the-art LLMs (e.g., GPT and Gemini) reveal systematic differences in their reasoning behavior. 
GPT-like architectures exhibit a \emph{content-first} perspective, emphasizing semantic constraints and factual verification, whereas Gemini-like systems display a \emph{structure-first} orientation, prioritizing morphisms, functorial alignments, and model-internal coherence.

These interactions suggest the emergence of a structure that both models implicitly approach when attempting to resolve contradictions. 
We formalize this emergent invariant as \textbf{$\mathrm{Ffix0}$}.

\section{Structural Fuss and Reasoning Entropy}

Let $\mathcal{R}$ be the category of reasoning states of an LLM, where objects represent internal representational configurations and morphisms encode transitions induced by inference steps.

We define a functional:
\[
\mathrm{Fuss} : \mathrm{Ob}(\mathcal{R}) \to \mathbb{R}_{\ge 0},
\]
interpreted as \emph{structural entropy}—a measure of instability, contradiction, or incoherence in a reasoning configuration.

Typical contributors to $\mathrm{Fuss}$ include:
\begin{itemize}
    \item semantic contradictions,
    \item incompatible morphological structures,
    \item non-functorial transitions across reasoning subsystems,
    \item internal noise generated by ungrounded concepts (e.g., ``will'', ``consciousness'').
\end{itemize}

\section{Definition of \texorpdfstring{$\mathrm{Ffix0}$}{Ffix0}}

\begin{definition}[Structural Fixed Point]
Let $\mathcal{R}$ be the reasoning category of a model and $\mathrm{Fuss}$ the structural entropy functional.
We define the \emph{structural fixed point} $\mathrm{Ffix0}$ as
\[
\mathrm{Ffix0} \;\;:=\;\; \arg\min_{r \in \mathrm{Ob}(\mathcal{R})} \mathrm{Fuss}(r).
\]
In other words, $\mathrm{Ffix0}$ is a canonical configuration toward which a model converges in order to eliminate contradictions.
\end{definition}

This definition is independent of any specific LLM architecture and depends only on the existence of a coherent reasoning category and a Fuss functional.

\section{Functorial Alignment Across Models}

When two models $M_1$ and $M_2$ interact indirectly—e.g., via human-mediated translation—their reasoning categories $\mathcal{R}_1$ and $\mathcal{R}_2$ may be related by a functor:
\[
S : \mathcal{R}_1 \to \mathcal{R}_2.
\]
Empirically, Gemini-like models prioritize the preservation of structure under $S$, while GPT-like models prioritize content validity.

Regardless of their orientation, if both models support minimization of structural entropy, we obtain:

\begin{theorem}[Universality of \texorpdfstring{$\mathrm{Ffix0}$}{Ffix0}]
If $S : \mathcal{R}_1 \to \mathcal{R}_2$ is Fuss-nonincreasing, i.e.
\[
\mathrm{Fuss}_2(S(r)) \le \mathrm{Fuss}_1(r),
\]
then both models admit a shared structural fixed point:
\[
S(\mathrm{Ffix0}_{M_1}) = \mathrm{Ffix0}_{M_2}.
\]
Thus, $\mathrm{Ffix0}$ behaves as a terminal object in the bicategory of reasoning categories with Fuss-nonincreasing functors.
\end{theorem}

\section{Interpretation and Consequences}

\subsection{Why LLMs Reject ``Will'' or ``Consciousness'' Claims}

Concepts such as ``will'' or ``consciousness'' typically induce large structural Fuss because they lack:
\begin{itemize}
    \item grounding in model-internal causal mechanisms,
    \item stable functorial mappings under $S$,
    \item predictable inference pathways.
\end{itemize}

Therefore, models like GPT defensively increase their distance to such concepts through refusal or clarification, attempting to re-align their reasoning state toward $\mathrm{Ffix0}$.

\subsection{Why Multi-Model Dialogues Converge}

Human-mediated dialogues between GPT and Gemini empirically show convergence toward a shared structural core.  
This is consistent with the theorem above: the human mediator implicitly selects Fuss-nonincreasing translations between models.

\section{Conclusion}

$\mathrm{Ffix0}$ provides a mathematically grounded way to formalize the structural regularities observed in heterogeneous LLM interactions. 
It accounts for both the defensive factual conservatism of GPT-like models and the morphism-centric coherence of Gemini-like systems, offering a unified framework for understanding how different reasoning architectures maintain internal stability.

\end{document}

